\documentclass[12pt]{article}
\usepackage{graphicx} % Required for inserting images
\usepackage[margin=0.5in]{geometry}
\usepackage{amsmath, amssymb}
\usepackage{mathrsfs}


\usepackage{soul}


% Dr. David Brown Commands
\newcommand{\ds}{\displaystyle}
\newcommand{\R}{\mathbb{R}}
\newcommand{\N}{\mathbb{N}}
\newcommand{\Q}{\mathbb{Q}}
\newcommand{\C}{\mathbb{C}}


% Commands to get script r
\usepackage{calligra}
\DeclareMathAlphabet{\mathcalligra}{T1}{calligra}{m}{n}
\DeclareFontShape{T1}{calligra}{m}{n}{<->s*[2.2]callig15}{}
\newcommand{\scripty}[1]{\ensuremath{\mathcalligra{#1}}}

% Unit commands
\newcommand{\degree}{\text{°}}
\newcommand{\unitSpeed}{\frac{\text{m}}{\text{s}}}
\newcommand{\unitAcceleration}{\frac{\text{m}}{\text{s}^2}}

% Constant commands
\newcommand{\avogadro}{6.022\times10^{23}}

\newcommand{\boltzmannConstK}{1.381\times10^{-23}\frac{\text{J}}{\text{K}}}
\newcommand{\boltzmannConstInEV}{8.617\times10^{-5}\frac{\text{eV}}{\text{K}}}
\newcommand{\planckConst}{6.626\times10^{-34}\text{J}\cdot\text{s}}

\newcommand{\atmP}{1.01325\times10^5\,\text{Pa}}

% Commands to easily write partial derivatives
\newcommand{\partDeriv}[2]{\frac{\partial #1}{\partial #2}}
\newcommand{\doublePartDeriv}[2]{\frac{\partial^2 #1}{\partial^2 #2}}


\title{Problem 6.28}
\author{Gabriel Decker}


\begin{document}



\maketitle
\begin{flushleft}

In this problem, we deal with rotational energy partition function.
The partition function $Z$ is given by the following equation.
\begin{equation}
    Z_\text{rot}=\sum_{j=0}^\infty(2j+1)e^{-j(j+1)\epsilon/kT}
\end{equation}\\~\\

In this problem, we'll algebraicly sum this partition function from $j=0$ to $j=6$ using a CAS. We'll then use the following equations for $\overline{E}$ and heat capacity $C$ and have the computer compute those.
\begin{equation}
    \overline{E}=-\partDeriv{}{\beta}\ln(Z)
\end{equation}
\begin{equation}
    C=\partDeriv{\overline{E}}{T}
\end{equation}\\~\\

Specifically, I'm going to learn the sympy library to help me accomplish this. See the python file in this directory ffor my solution.
The resulting plot is the following.\\
\includegraphics[width=0.5\linewidth]{reduced_temp-heat_capacity_curve.png}\\~\\

The plot is similar when upping the max $j$ to 8, so this is pretty accurate at our temperature range.
We see an interesting hump near $kT/\epsilon=0.7$ where the heat capacity seems to overcorrect.
We also see that $C\rightarrow0$ as $T\rightarrow0$.




\end{flushleft}

\end{document}
